%
% teil2.tex -- Beispiel-File für teil2 
%
% (c) 2020 Prof Dr Andreas Müller, Hochschule Rapperswil
%
% !TEX root = ../../buch.tex
% !TEX encoding = UTF-8
%
\section{Der Weierstraß-Produktsatz}
\rhead{Weierstraß-Produktsatz}

\subsection{Elementarfaktoren}

Das zentrale technische Hilfsmittel sind die Weierstraß-Elementarfaktoren.
Für $p = 0, 1, 2, \ldots$ definiert man:
\begin{equation}
E_0(z) = 1 - z,
\qquad
E_p(z) = (1-z) \cdot \exp\!\left(z + \frac{z^2}{2} + \cdots + \frac{z^p}{p}\right)
\quad \text{für } p \ge 1.
\end{equation}
Der Exponentialterm ist so konstruiert, dass er die Divergenz des Faktors $(1-z)$
für wachsende Argumente kompensiert.
Genauer gilt die Abschätzung
\begin{equation}
|1 - E_p(z)| \le |z|^{p+1}
\quad \text{für } |z| \le 1.
\end{equation}
Diese Ungleichung ist der Schlüssel für den Konvergenzbeweis des Produktsatzes;
für einen vollständigen Beweis sei auf \cite{produkt:conway}
oder \cite{produkt:steinShakarchi} verwiesen.

\subsection{Formulierung des Satzes}

Weierstraß bewies den folgenden fundamentalen Satz \cite{produkt:steinShakarchi}:

\begin{satz}[Weierstraß, 1876]
Sei $f$ eine ganze Funktion mit einer Nullstelle der Ordnung $m \ge 0$ bei $z = 0$
und weiteren Nullstellen $a_1, a_2, \ldots$ (mit Vielfachheiten,
nach wachsendem Betrag geordnet).
Dann existiert eine ganze Funktion $g(z)$ und nicht-negative ganze Zahlen $p_n$ mit
\begin{equation}
f(z)
=
z^m \cdot e^{g(z)} \cdot \prod_{n=1}^{\infty} E_{p_n}\!\left(\frac{z}{a_n}\right).
\end{equation}
Das Produkt konvergiert gleichmäßig auf Kompakta.
Für ganze Funktionen endlicher Ordnung $\varrho$ genügt die konstante Wahl
$p_n = \lfloor \varrho \rfloor$.
\end{satz}

\subsection{Rückkehr zum Sinus}

Die Sinusfunktion $\sin(\pi z)$ ist eine ganze Funktion der Ordnung $\varrho = 1$,
da $|\sin(\pi z)| \le e^{\pi|z|}$ und die Nullstellen $a_n = \pm n$ linear wachsen.
Man wählt daher $p_n = 1$, also $E_1(z) = (1-z)e^z$.
Wegen der Symmetrie $\sin(\pi(-z)) = -\sin(\pi z)$ heben sich die Exponentialterme
der Paare $+n$ und $-n$ gegenseitig auf -- denn $e^{z/n} \cdot e^{-z/n} = 1$ --,
und man erhält:
\begin{equation}
\sin(\pi z)
=
\pi z \cdot e^{g(z)} \cdot \prod_{n=1}^{\infty} \left(1 - \frac{z^2}{n^2}\right).
\end{equation}
Der Hadamard-Faktorisierungssatz -- eine Schärfung des Weierstraß-Satzes
für Funktionen endlicher Ordnung -- impliziert, dass $g(z)$ ein Polynom vom Grad
höchstens $\lfloor \varrho \rfloor = 1$ sein muss, also $g(z) = Az + B$.
Aus $\sin(\pi z)/(\pi z) \to 1$ für $z \to 0$ folgt $B = 0$,
und die Symmetrie, dass $\sin(\pi z)$ ungerade ist, ergibt $A = 0$.
Damit gilt Eulers Formel rigoros:
\begin{equation}
\sin(\pi z)
=
\pi z \cdot \prod_{n=1}^{\infty} \left(1 - \frac{z^2}{n^2}\right).
\end{equation}
